% Options for packages loaded elsewhere
% Options for packages loaded elsewhere
\PassOptionsToPackage{unicode}{hyperref}
\PassOptionsToPackage{hyphens}{url}
\PassOptionsToPackage{dvipsnames,svgnames,x11names}{xcolor}
%
\documentclass[
  letterpaper,
  DIV=11,
  numbers=noendperiod]{scrartcl}
\usepackage{xcolor}
\usepackage{amsmath,amssymb}
\setcounter{secnumdepth}{-\maxdimen} % remove section numbering
\usepackage{iftex}
\ifPDFTeX
  \usepackage[T1]{fontenc}
  \usepackage[utf8]{inputenc}
  \usepackage{textcomp} % provide euro and other symbols
\else % if luatex or xetex
  \usepackage{unicode-math} % this also loads fontspec
  \defaultfontfeatures{Scale=MatchLowercase}
  \defaultfontfeatures[\rmfamily]{Ligatures=TeX,Scale=1}
\fi
\usepackage{lmodern}
\ifPDFTeX\else
  % xetex/luatex font selection
\fi
% Use upquote if available, for straight quotes in verbatim environments
\IfFileExists{upquote.sty}{\usepackage{upquote}}{}
\IfFileExists{microtype.sty}{% use microtype if available
  \usepackage[]{microtype}
  \UseMicrotypeSet[protrusion]{basicmath} % disable protrusion for tt fonts
}{}
\makeatletter
\@ifundefined{KOMAClassName}{% if non-KOMA class
  \IfFileExists{parskip.sty}{%
    \usepackage{parskip}
  }{% else
    \setlength{\parindent}{0pt}
    \setlength{\parskip}{6pt plus 2pt minus 1pt}}
}{% if KOMA class
  \KOMAoptions{parskip=half}}
\makeatother
% Make \paragraph and \subparagraph free-standing
\makeatletter
\ifx\paragraph\undefined\else
  \let\oldparagraph\paragraph
  \renewcommand{\paragraph}{
    \@ifstar
      \xxxParagraphStar
      \xxxParagraphNoStar
  }
  \newcommand{\xxxParagraphStar}[1]{\oldparagraph*{#1}\mbox{}}
  \newcommand{\xxxParagraphNoStar}[1]{\oldparagraph{#1}\mbox{}}
\fi
\ifx\subparagraph\undefined\else
  \let\oldsubparagraph\subparagraph
  \renewcommand{\subparagraph}{
    \@ifstar
      \xxxSubParagraphStar
      \xxxSubParagraphNoStar
  }
  \newcommand{\xxxSubParagraphStar}[1]{\oldsubparagraph*{#1}\mbox{}}
  \newcommand{\xxxSubParagraphNoStar}[1]{\oldsubparagraph{#1}\mbox{}}
\fi
\makeatother


\usepackage{longtable,booktabs,array}
\newcounter{none} % for unnumbered tables
\usepackage{calc} % for calculating minipage widths
% Correct order of tables after \paragraph or \subparagraph
\usepackage{etoolbox}
\makeatletter
\patchcmd\longtable{\par}{\if@noskipsec\mbox{}\fi\par}{}{}
\makeatother
% Allow footnotes in longtable head/foot
\IfFileExists{footnotehyper.sty}{\usepackage{footnotehyper}}{\usepackage{footnote}}
\makesavenoteenv{longtable}
\usepackage{graphicx}
\makeatletter
\newsavebox\pandoc@box
\newcommand*\pandocbounded[1]{% scales image to fit in text height/width
  \sbox\pandoc@box{#1}%
  \Gscale@div\@tempa{\textheight}{\dimexpr\ht\pandoc@box+\dp\pandoc@box\relax}%
  \Gscale@div\@tempb{\linewidth}{\wd\pandoc@box}%
  \ifdim\@tempb\p@<\@tempa\p@\let\@tempa\@tempb\fi% select the smaller of both
  \ifdim\@tempa\p@<\p@\scalebox{\@tempa}{\usebox\pandoc@box}%
  \else\usebox{\pandoc@box}%
  \fi%
}
% Set default figure placement to htbp
\def\fps@figure{htbp}
\makeatother





\setlength{\emergencystretch}{3em} % prevent overfull lines

\providecommand{\tightlist}{%
  \setlength{\itemsep}{0pt}\setlength{\parskip}{0pt}}



 


\KOMAoption{captions}{tableheading}
\makeatletter
\@ifpackageloaded{caption}{}{\usepackage{caption}}
\AtBeginDocument{%
\ifdefined\contentsname
  \renewcommand*\contentsname{Table of contents}
\else
  \newcommand\contentsname{Table of contents}
\fi
\ifdefined\listfigurename
  \renewcommand*\listfigurename{List of Figures}
\else
  \newcommand\listfigurename{List of Figures}
\fi
\ifdefined\listtablename
  \renewcommand*\listtablename{List of Tables}
\else
  \newcommand\listtablename{List of Tables}
\fi
\ifdefined\figurename
  \renewcommand*\figurename{Figure}
\else
  \newcommand\figurename{Figure}
\fi
\ifdefined\tablename
  \renewcommand*\tablename{Table}
\else
  \newcommand\tablename{Table}
\fi
}
\@ifpackageloaded{float}{}{\usepackage{float}}
\floatstyle{ruled}
\@ifundefined{c@chapter}{\newfloat{codelisting}{h}{lop}}{\newfloat{codelisting}{h}{lop}[chapter]}
\floatname{codelisting}{Listing}
\newcommand*\listoflistings{\listof{codelisting}{List of Listings}}
\makeatother
\makeatletter
\makeatother
\makeatletter
\@ifpackageloaded{caption}{}{\usepackage{caption}}
\@ifpackageloaded{subcaption}{}{\usepackage{subcaption}}
\makeatother
\usepackage{bookmark}
\IfFileExists{xurl.sty}{\usepackage{xurl}}{} % add URL line breaks if available
\urlstyle{same}
% fallback for those not using the hyperref driver hyperxmp:
\makeatletter
\@ifundefined{xmpquote}{\newcommand{\xmpquote}[1]{#1}}{}
\makeatother
\hypersetup{
  pdftitle={Reproducibility Report for Study by Salomons et al.~(2026{,} Animal Behaviour)},
  pdfauthor={Samantha Kefer (skefer@ucsd.edu)},
  colorlinks=true,
  linkcolor={blue},
  filecolor={Maroon},
  citecolor={Blue},
  urlcolor={Blue},
  pdfcreator={LaTeX via pandoc}}


\title{Reproducibility Report for Study by Salomons et al.~(2026, Animal
Behaviour)}
\author{Samantha Kefer (skefer@ucsd.edu)}
\date{2026-10-25}
\begin{document}
\maketitle

\renewcommand*\contentsname{Table of contents}
{
\setcounter{tocdepth}{3}
\tableofcontents
}

\subsection{Introduction}\label{introduction}

{[}No abstract is needed.{]} Each reproducibility project will have a
straightforward, no frills report of the study and reproducibility
results. These reports will be publicly available as supplementary
material for the aggregate report(s) of the project as a whole. Also, to
maximize project integrity, the intro and methods will be written and
critiqued in advance of your attempt to reproduce the results.
Introductions can be just 1-2 paragraphs clarifying the main idea of the
original study, the target finding for your reproducibility attempt, and
any other essential information. It will NOT have a literature review --
that is in the original publication. You can write both the introduction
and the methods in past tense.

\subsubsection{Justification for choice of
study}\label{justification-for-choice-of-study}

The paper by Salomons et al.~explains the cognitive development of
assistance puppies (aged 8 to 20 weeks) as measured through performance
on a biweekly battery of 9 tasks: marker gesture, working memory,
working memory with distraction, auditory discrimination, odor
discrimination, cylinder (inhibitory control and reversal learning),
unsolvable, causal reasoning, and momentary pointing. The results
support the multiple intelligences hypothesis and suggest the emergence
of distinct cognitive abilities, with clear developmental breakpoints in
marker gesture comprehension, inhibitory control, reversal learning, and
eye contact occurring before 16 weeks, coinciding with the period of
rapid brain development. The goal of this project is to answer if
different cognitive skills develop at different rates in assistance
puppies. The authors, however, chose to model each task separately and
compare developmental rate visually rather than test comparisons
directly. Using data exclusively from the 91 longitudinal subjects, I
will reproduce Figure 2 and directly test whether age-related
improvement differs across tasks using an age × task interaction. Next,
I will simulate additional data from the fitted model to test whether a
study of this size (N=91) can reliably detect these differences or is a
larger sample (N=1000) is needed for sufficient power.

I am currently working to develop a battery of tests that can be
administered to large numbers of dogs, in the hopes of identifying
exceptional performers and determining whether strong performance on one
task may generalize across different measures of the same ability and/or
across cognitive domains. Knowing when different skills emerge will
inform when puppies should be tested, and between-task correlations will
suggest whether performance may generalize across domains. With the help
of citizen science, this method may ultimately identify exceptional
dogs.

\subsubsection{Anticipated challenges}\label{anticipated-challenges}

As explained above, my plan for this project includes (1) reanalyzing
existing data from the 91 longitudinal puppies and (2) simulating
additional subjects. My plan for this project will include 6 steps:

\begin{enumerate}
\def\labelenumi{\arabic{enumi}.}
\tightlist
\item
  Reproduce Figure 2 with LOESS curve
\item
  Reproduce per-task models correct \textasciitilde{} age + (1
  \textbar{} puppy\_ID) and correct \textasciitilde{} age + experience +
  (1 \textbar{} puppy\_ID)
\item
  Confirmatory mixed-effects logistic model, correct \textasciitilde{}
  age * task + (1 \textbar{} puppy\_ID) + (1 \textbar{} puppy\_ID:task)

  \begin{itemize}
  \tightlist
  \item
    Key result is age × task interaction (only including data from the
    91 longitudinal subjects, so age and experience will confound)
  \item
    Significant interaction supports multiple intelligences hypothesis
  \end{itemize}
\item
  Calculate effect sizes per task and between tasks (for how much faster
  one skill develops relative to another)
\item
  Data simulation using fitted model (step 3) for N=91 vs N=1000
\item
  Plot simulated data with LOESS curves like Figure 2
\end{enumerate}

My reproduction will be successful if curves match Figure 2, betas are
within +/- 0.01 of the reported, and a find the same best-AIC model per
task. My simulation will be successful if slopes show little bias, I can
get \textasciitilde95\% CI coverage, and \textasciitilde80\% power for
interaction, ideally at N=91.

I anticipate running into a few challenges when attempting to reproduce
these results. First, age and experience are highly intertwined since
they increase together for all longitudinal subjects. In addition, many
of the dogs are siblings, half-siblings, cousins, etc. and this
relatedness is not included in the dataset and may be a confound.
Finally, of the 11 measures, 10 are binary responses (score out of 4
trials) and one is average looking time, thus I am considering removing
the last measure from my simulation for ease. The biggest challenge,
however, will be actually completing this project given my minimal
programming knowledge and experience.

\subsubsection{Learning Goals}\label{learning-goals}

\begin{itemize}
\tightlist
\item
  learn to fit and interpret mixed-effects models for the data
\item
  learn how to write a simulation function that generates correlated
  data across all 11 measures for the same subject
\item
  learn how to use simulation to test how sample size and testing
  schedule affect power (directly applies to planning my future project)
\end{itemize}

\subsubsection{Links}\label{links}

Project repository (on Github):
https://github.com/samkefer/salomons2026/tree/main

Original paper (as hosted in your repo):
https://github.com/samkefer/salomons2026/blob/main/original\_paper/original\_paper

\subsection{Methods}\label{methods}

\subsubsection{Description of the steps required to reproduce the
results}\label{description-of-the-steps-required-to-reproduce-the-results}

Please describe all the steps necessary to reproduce the key result(s)
of this study.

\subsubsection{Differences from original
study}\label{differences-from-original-study}

Explicitly describe known differences in the analysis pipeline between
the original paper and yours (e.g., computing environment). The goal, of
course, is to minimize those differences, but differences may occur.
Also, note whether such differences are anticipated to influence your
ability to reproduce the original results.

\subsection{Project Progress Check 1}\label{project-progress-check-1}

\subsubsection{Measure of success}\label{measure-of-success}

Please describe the outcome measure for the success or failure of your
reproduction and how this outcome will be computed.

\subsubsection{Pipeline progress}\label{pipeline-progress}

Earlier in this report, you described the steps necessary to reproduce
the key result(s) of this study. Please describe your progress on each
of these steps (e.g., data preprocessing, model fitting, model
evaluation).

\subsection{Results}\label{results}

\subsubsection{Data preparation}\label{data-preparation}

Data preparation following the analysis plan.

\subsubsection{Key analysis}\label{key-analysis}

The analyses as specified in the analysis plan.

\emph{Side-by-side graph with original graph is ideal here}

\subsubsection{Exploratory analyses}\label{exploratory-analyses}

Any follow-up analyses desired (not required).

\subsection{Discussion}\label{discussion}

\subsubsection{Summary of Reproduction
Attempt}\label{summary-of-reproduction-attempt}

Open the discussion section with a paragraph summarizing the primary
result from the key analysis and assess whether you successfully
reproduced it, partially reproduced it, or failed to reproduce it.

\subsubsection{Commentary}\label{commentary}

Add open-ended commentary (if any) reflecting (a) insights from
follow-up exploratory analysis of the dataset, (b) assessment of the
meaning of the successful or unsuccessful reproducibility attempt -
e.g., for a failure to reproduce the original findings, are the
differences between original and present analyses ones that definitely,
plausibly, or are unlikely to have been moderators of the result, and
(c) discussion of any objections or challenges raised by the current and
original authors about the reproducibility attempt (if you contacted
them). None of these need to be long.




\end{document}
